% Architecture: the parts, where each runs, trust boundaries and deployment topologies, a
% request's path through the server, sign-in, the protection of patient data in brief, the
% gateway, data flow, readers and data models, extension points, contracts, technology, the code
% map and the status. Written for IT security reviewers, integrators and contributors.
% Source: docs/architecture.md; the code is the authority (medorama/src/server.rs router, guard,
% privacy; signin.rs; gateway.rs; viewer.rs; jobs.rs; sources/; libs/engine).
\chapter{Architecture}
\label{chap:architecture}

This chapter describes how Medorama is built: its parts, where each one runs, how a request
travels through the server, how data come in, and where the system can be adapted. It is written
for institutional IT security reviewers, for integrators who deploy Medorama, and for
contributors, and it can be read on its own. How patients' identities are protected is explained
in \cref{chap:protecting}, and the security controls are reviewed, with their code and tests, in
\cref{chap:security}. Together, these three chapters form the part of the manual meant for a
review before Medorama is deployed on patient data.

\section{Overview}
\label{sec:architecture-overview}

Medorama is a research system that gathers each patient's data from many sources (admissions,
laboratory results, vital signs, medications, notes, ECGs, imaging and bedside-monitor
recordings) into one \emph{catalog}, Medorama's own database. It shows each patient's journey on
one page, lets researchers build cohorts across all sources, and opens recordings and images in
its own viewers. It is one server program that an institution runs on its own hosts, and users
reach it with a web browser. Every request passes Medorama's sign-in, and every answer to a user
who may not see identifiers leaves through the \emph{privacy gate}, which replaces identifiers
and moves dates back (\cref{fig:architecture-glance}).

\screenshot{architecture.pdf}{Medorama at a glance: users reach one server through its sign-in;
the server reads the catalog, the data sources and connected databases, runs the viewers and
plug-ins, and sends every answer of a session in Privacy Mode through the privacy
gate.}{fig:architecture-glance}

\section{Design principles}
\label{sec:architecture-principles}

\begin{enumerate}
\item \textbf{Separate programs, joined by written contracts.} Each part builds, runs and is
  tested on its own. Parts meet only through a documented contract: a file format, an HTTP API,
  a command-line protocol or a folder layout (\cref{sec:architecture-contracts}).
\item \textbf{One way in.} When Medorama fronts the viewers, every request passes its sign-in,
  roles and rules. The viewers stay simple and read only what Medorama hands them.
\item \textbf{One way out.} The catalog keeps every identifier as imported. In Privacy Mode,
  every answer the server sends (JSON, files, exports, the waveform viewer's answers) passes the
  privacy gate, and the server's last layer refuses any answer that the gate did not clear, so a
  new endpoint cannot forget the rule.
\item \textbf{Sources are only read.} No part of Medorama writes into a data folder or a
  connected database. Derived data (indexes, results) go to their own folders or to the catalog.
\item \textbf{The server decides; the page draws.} What a user may see, and how data are
  matched, merged and searched, is decided on the server. The pages draw what they are sent,
  with no build step and no third-party code.
\item \textbf{Open formats inside, licensed formats outside.} The repository reads open formats
  (WFDB, EDF, DICOM, GE MUSE XML, tables). A format under a vendor's license is read by an
  \emph{add-on}, a separate program that an administrator installs and that is distributed under
  its own terms. Nothing licensed is compiled into Medorama or a viewer.
\item \textbf{Recordings are read identically three times.} The VitalStream viewer has three
  independent implementations (Rust, C++ and Python) that write the same index byte for byte and
  answer the same API. They are checked against each other, so that a misread recording does not
  go unnoticed.
\end{enumerate}

\section{The parts}
\label{sec:architecture-parts}

\Cref{tab:architecture-parts} lists the parts of Medorama, where each one lives in the
repository, and how it runs.

\begin{xltabular}{\linewidth}{@{}P{3.0cm}P{2.9cm}L@{}}
\caption{The parts of Medorama.}
\label{tab:architecture-parts}\\
\toprule
\tabhead Part & \tabhead Folder; language & \tabhead Runs as \\
\midrule
\endfirsthead
\toprule
\tabhead Part & \tabhead Folder; language & \tabhead Runs as \\
\midrule
\endhead
\bottomrule
\endfoot
Medorama server & \file{medorama/}; Rust & One program, \cmd{medorama serve}, on the
  institution's host; also the \cmd{medorama} command line for administration \\
Medorama's web page & \file{medorama/web/}; HTML, CSS, JavaScript & Static files embedded in the
  server and run in the user's browser, with light and dark themes \\
VitalStream (waveforms: ICU monitoring, Holter, EEG, sleep) & \file{viewers/vitalstream/}; Rust,
  C++, Python, and a shared page & Alone, a local program with its own page. Inside Medorama, a
  child process on a private loopback port, behind the gateway \\
ECG viewer (12-lead ECGs on standard paper) & \file{viewers/ecg/}; JavaScript, and
  \cmd{muse-py} in Python & A web page, opened from disk, served by \cmd{muse-py serve}, or
  served by Medorama as a tab \\
PACS viewer (DICOM studies) & \file{viewers/pacs/}; JavaScript, and \cmd{dicom-py} in Python &
  A web page, like the ECG viewer (\cmd{dicom-py serve} alone) \\
Engine & \file{libs/engine/} (\code{recording-engine}), \file{libs/signal-core/}; Rust & A
  library used by Medorama and the Rust viewer: readers, the index, queries, the adapter
  protocol. WFDB is read through wfdb-x \\
Add-ons (adapters) & Outside this repository; any language & Programs in the adapter folder, run
  as child processes; for example stp-adapter for BedMasterEx \file{.Stp} files (private,
  available to institutions under agreement) \\
Plug-in actions & An administrator's \code{actions\_dirs}; any language & Programs run as child
  processes by the job queue (\cref{chap:actions}) \\
Catalog database & The catalog's home folder, or a PostgreSQL server & Medorama's own data:
  users, sessions, the audit log, settings, patients and identifiers, pseudonyms, cohorts, the
  job queue \\
Indexes and results & \code{index\_dir}, \code{results\_dir}; files & Waveform indexes (the
  viewers' format) and the outputs of actions \\
Data sources & \code{data\_roots}; connected databases & Exported tables, recordings, ECGs,
  DICOM studies and documents, read where they lie; OMOP or other databases read through
  views \\
\end{xltabular}

The viewers are programs of their own. Each is useful alone and never depends on Medorama.
Medorama opens them as tabs, at an address that names the file to show
(\cref{chap:viewers-in-medorama}).

\section{Trust boundaries and deployment}
\label{sec:architecture-boundaries}

\screenshot{system-context.pdf}{The system context and its trust boundaries: the browser, the
optional authenticating proxy, Medorama's server with its viewer and plug-ins, the data
sources, the catalog's database and the connected databases.}{fig:architecture-context}

The browser and the network are untrusted: every request is authenticated and authorized on the
server, and the page decides nothing about access. The server host is trusted: one
operating-system account, the catalog's owner, runs Medorama, its viewer and its plug-ins. The
zones and how far each is trusted are tabulated in \cref{tab:security-zones}.

\subsection{What listens where}
\label{sec:architecture-listen}

\begin{xltabular}{\linewidth}{@{}P{3.8cm}P{4.4cm}L@{}}
\caption{The programs that listen, and who reaches them.}
\label{tab:architecture-listen}\\
\toprule
\tabhead Program & \tabhead Listens on (default) & \tabhead Reached by \\
\midrule
\endfirsthead
\toprule
\tabhead Program & \tabhead Listens on (default) & \tabhead Reached by \\
\midrule
\endhead
\bottomrule
\endfoot
\cmd{medorama serve} & 127.0.0.1:8090 (\code{host}, \code{port}) & Browsers, through a TLS
  reverse proxy or an SSH tunnel \\
VitalStream under Medorama & 127.0.0.1, on a free port chosen at start & Medorama's gateway only,
  with requests signed by a key only Medorama knows \\
VitalStream alone & 127.0.0.1:8050 (Python), 8060 (Rust), 8070 (C++) & A browser on the same
  computer, or through an SSH tunnel \\
\cmd{muse-py serve}, \cmd{dicom-py serve} & 127.0.0.1:8080, 127.0.0.1:8081 & The same \\
Adapters and plug-ins & Nothing & Medorama and the viewers start them and talk to them over
  pipes and files \\
\end{xltabular}

Medorama listens on the loopback address by default, so that nothing on the network can reach
it until an administrator puts a TLS proxy in front of it. If it is bound to another address
without \code{secure\_cookies = true}, it prints a warning at start. It opens outgoing
connections only to its PostgreSQL catalog and to connected databases, always over verified TLS
when they are on another host, and it sends nothing anywhere else.

\subsection{What is read and what is written}
\label{sec:architecture-data}

\begin{xltabular}{\linewidth}{@{}P{5.4cm}P{3.8cm}L@{}}
\caption{Where data live, and who writes them.}
\label{tab:architecture-data}\\
\toprule
\tabhead Data & \tabhead Lives in & \tabhead Written by \\
\midrule
\endfirsthead
\toprule
\tabhead Data & \tabhead Lives in & \tabhead Written by \\
\midrule
\endhead
\bottomrule
\endfoot
Recordings, ECGs, images, documents, exported tables & Their own folders, inside
  \code{data\_roots} & Nobody: they are only read \\
Connected databases & The institution's database servers & Nobody: they are read through
  read-only sessions \\
Patients, encounters, entries, identifiers, pseudonyms, users, sessions, jobs, the audit log &
  The catalog database & Medorama \\
Waveform indexes & \code{index\_dir} (which standalone viewers' \code{cache\_dir} may share) &
  The engine, the viewers, adapters \\
Results of actions & \code{results\_dir} & Plug-ins, through the job queue \\
\end{xltabular}

\subsection{Deployment topologies}
\label{sec:architecture-topologies}

\begin{itemize}
\item \textbf{One server with SQLite.} This is the default. \cmd{medorama init} creates a home
  folder with \file{medorama.toml} and the SQLite database (readable by its owner only, and
  optionally encrypted with SQLCipher). One \cmd{medorama serve} process answers requests and
  runs jobs in worker threads. This suits a team or a department.
\item \textbf{Several servers with PostgreSQL.} Identical servers share one catalog in
  PostgreSQL (\code{database\_url}), with the same \code{data\_roots} and with the index and
  results folders on shared storage. Any server answers any request. Sessions, cohorts and the
  job queue are in the database, and work that must happen once takes an advisory lock. A load
  balancer or reverse proxy spreads the requests (\cref{chap:administration}).
\item \textbf{Behind a TLS reverse proxy.} nginx, Apache httpd or a load balancer terminates
  HTTPS and passes requests to Medorama on 127.0.0.1, with \code{secure\_cookies = true}. Users
  sign in with Medorama's passwords (\cref{chap:deployment}).
\item \textbf{Behind the institution's sign-in.} An authenticating proxy (oauth2-proxy, Apache
  httpd with mod\_auth\_openidc, Shibboleth or Kerberos, nginx with \code{auth\_request}) signs
  users in with single sign-on and multi-factor authentication, and names them to Medorama in a
  header (\code{sign\_in = "proxy"}; \cref{chap:sign-in}).
\item \textbf{Viewers alone on a workstation.} A researcher runs a VitalStream viewer, or opens
  the ECG or PACS viewer's page, on their own computer, with no Medorama and no catalog. The
  viewer reads files under a root folder fixed when it was built, and may ask for a password
  (\cref{sec:sign-in-viewers}).
\end{itemize}

\subsection{Running separately and together}
\label{sec:architecture-together}

\textbf{Separately}, each viewer runs on its own: the VitalStream viewers on their ports, with
their own page and optional password (\file{viewers/vitalstream/install.sh}), and the ECG and
PACS viewers as pages opened from disk or served by \cmd{muse-py serve} and \cmd{dicom-py serve}.
Medorama runs on its own port and may link to a VitalStream viewer that runs elsewhere
(\code{viewer\_url}).

\textbf{Together} (\cmd{tools/install-medorama.sh -{}-viewer cpp|rs|py|none}), Medorama is
installed with one VitalStream viewer, and the ECG and PACS viewers are built into Medorama.
Medorama's supervisor (\file{viewer.rs}) starts the VitalStream viewer on first use, on
127.0.0.1 and a free port, with a settings file and a private password file that Medorama
writes. It restarts the viewer if its process ends, or if a request cannot reach it and it then
fails a health check, and stops it when Medorama stops. A request to a running viewer waits for
no other request and costs no check, so a viewer busy with a large recording is never taken for
a dead one. Add-ons are
installed into the adapter folder that Medorama and the viewers share
(\cmd{medorama addons install PROGRAM}).

\section{A request's path through the server}
\label{sec:architecture-request}

\screenshot{components.pdf}{Inside the server: the path of one request, through the guard, the
sign-in and the privacy layer to its handler, and back out through the privacy
gate.}{fig:architecture-components}

Every HTTP request passes the same layers, which are axum middleware (\file{server.rs},
\code{router}):
\begin{enumerate}
\item \textbf{The guard} (\code{guard}). The server refuses any API request that changes
  something unless it carries the header \code{X-Medorama: 1}, which a page on another site
  cannot add. On the way out, the guard sets security headers on every answer: a Content
  Security Policy, framing rules, \code{nosniff}, \code{no-referrer}, and \code{no-store} on the
  API.
\item \textbf{Sign-in} (\file{signin.rs}, \code{identify}). The privacy layer asks the sign-in
  module, once per request, who the request comes from: the account of a session cookie, or the
  account that an authenticating proxy names.
\item \textbf{The privacy layer, inbound} (\code{privacy}). The server decides whether the
  session is in Privacy Mode. In Privacy Mode, the user is replaced by a copy without the PHI
  permission.
\item \textbf{The handler.} The route's handler receives the user as \code{Auth(User)} (HTTP 401
  without one), checks the role and the permission it needs, and reads the catalog through the
  core modules. A handler that knows its data de-identifies them itself (files, table cells,
  notes).
\item \textbf{The privacy layer, outbound: the gate.} For a session in Privacy Mode, every JSON
  answer passes the privacy gate, a file passes only if its handler cleared it, and any other
  answer carrying data is refused. The answer then leaves through the guard.
\end{enumerate}
The handlers are thin: the rules live in modules that the command line, the background workers
and the tests also use.

\section{Sign-in}
\label{sec:architecture-sign-in}

How users sign in is decided in one module, \file{signin.rs}, and chosen by the setting
\code{sign\_in}:
\begin{itemize}
\item \code{password} (the default) uses Medorama's own accounts, with Argon2id password hashes,
  throttling, and session cookies holding a random token that is stored hashed;
\item \code{proxy} lets the institution's authenticating reverse proxy sign the user in (single
  sign-on through OpenID Connect or SAML, Kerberos, multi-factor authentication). The proxy names
  the user in a header, which Medorama believes only on connections from the proxy's addresses.
\end{itemize}
With either method, accounts, roles and the PHI permission live in Medorama, and every handler
receives the signed-in user in the same form (\cref{fig:sign-in-seam}). Supporting another
method means changing the sign-in module only. The settings, proxy examples and the code changes
for a new method are in \cref{chap:sign-in}.

\section{Protecting patient data}
\label{sec:architecture-privacy}

Identifiers (names, medical record numbers, dates of birth, account numbers, imported IDs) are
kept only in the catalog database, in their own tables. What leaves the server is decided at the
moment of sending. A session is in \emph{Privacy Mode} unless its user has the PHI permission and
has switched Privacy Mode off. In Privacy Mode, the privacy gate (\file{gate.rs}, \file{deid.rs},
\file{dicom.rs}, \file{pseudonyms.rs}):
\begin{itemize}
\item replaces identifiers with placeholders (\code{[Patient Name]}, \code{[MRN]}) and random
  pseudonyms (\code{MO-XXXX-XXXX});
\item moves every date back by the catalog's offset plus each patient's own number of days, so
  that a patient's intervals are kept;
\item de-identifies files at the source (WFDB headers, GE MUSE XML, DICOM headers following
  PS3.15's Basic Profile, EDF and BDF headers and annotations), passes pictures, videos and PDF
  documents with a warning, and refuses any other file;
\item refuses searches of identifiers;
\item applies to exports, to the VitalStream viewer through the gateway, and to connected
  databases' data; plug-ins receive pseudonyms only.
\end{itemize}
The full explanation is \cref{chap:protecting}, and the controls, their code and their tests are
in \cref{sec:security-privacy}.

\section{The gateway}
\label{sec:architecture-gateway}

When Medorama runs a VitalStream viewer, the browser never talks to the viewer directly.
Medorama passes the viewer's requests through \file{/viewer/} (\file{gateway.rs}) and decides
each one.
\begin{itemize}
\item \textbf{Default deny.} Only the listed endpoints of the viewer's API (version 1) pass.
  Anything else, including endpoints that a newer viewer may add, is refused until it is given a
  rule.
\item \textbf{Medorama's sign-in.} A signed-in Medorama user is required. The viewer runs with a
  private password that Medorama generates; Medorama signs each request it passes on and never
  forwards the user's own cookie. Users never see the viewer's sign-in, and the viewer's port is
  of no use to anyone else on the host.
\item \textbf{Rules per user}, as \cref{tab:architecture-gateway} shows.
\item \textbf{Audit.} Each opened recording is written to the audit log (\code{view}).
\item \textbf{Paths} are compared after symbolic links are resolved, as the viewers do.
\item \textbf{Settings.} The viewer starts with a settings file that Medorama writes (its index
  and adapter folders among them), so that a user's own \file{vitalstream.toml} cannot conflict
  with the gateway's rules.
\end{itemize}

\begin{xltabular}{\linewidth}{@{}P{4.6cm}P{4.4cm}L@{}}
\caption{The gateway's rules for each viewer request.}
\label{tab:architecture-gateway}\\
\toprule
\tabhead Viewer request & \tabhead Administrators & \tabhead Analysts and viewers \\
\midrule
\endfirsthead
\toprule
\tabhead Viewer request & \tabhead Administrators & \tabhead Analysts and viewers \\
\midrule
\endhead
\bottomrule
\endfoot
The Open dialog (\code{browse}, \code{complete}, places and the start folder in \code{config}) &
  Folders inside \code{data\_roots} (everything when it is empty) & Only folders that hold
  recordings in the catalog, showing only those recordings \\
\code{scan} a folder and \code{open} a recording & Any recording inside \code{data\_roots} & Only
  recordings in the catalog \\
\code{cancel} an index being built; \code{datasets/\{key\}\dots} (windows, vital signs, markers)
  & Only datasets they opened through the gateway & The same \\
In Privacy Mode (anyone) & Browsing and scanning refused; recordings opened by token
  (\code{mo-rec:<n>}); times translated both ways; paths and names replaced; annotation texts
  treated like notes; the viewer's error texts (which name files) replaced by a fixed sentence,
  and the records of a folder it left out named by number & The same \\
\end{xltabular}

\begin{itemize}
\item \textbf{Changes are JSON.} A request that changes something (\code{open}, \code{cancel})
  must be sent as \code{application/json} with \code{X-Medorama: 1}, which another site's page
  cannot send unasked.
\item \textbf{The page's policy.} The viewer's page is served with a Content Security Policy of
  its own scripts and styles only (\code{default-src 'self'; script-src 'self'; img-src 'self'
  data:; object-src 'none'; base-uri 'none'; frame-ancestors 'self'}).
\end{itemize}

The ECG and PACS viewers need no gateway. They are pages that Medorama serves, and they read only
an entry's files, through \code{patients/\{id\}/file/\{item\}}, which Medorama checks like every
other request (\cref{chap:http-api}).

\section{Data flow}
\label{sec:architecture-flow}

\paragraph{Importing databases.} Data enter the catalog as databases
(\cref{fig:importing-flow}). An administrator imports a table (CSV, Excel, OpenDocument or
Parquet) exported from a source system, from the page or the command line. Medorama recognizes
its columns (patient ID, MRN, times, values, file), keeps the column names as imported, matches
each row to a patient and an encounter, and stores the identifiers in the identifier tables. Rows
that it cannot place wait in the \ui{Review} tab. Files are linked where they lie, by the rows
that name their patient, and only inside \code{data\_roots}. Known exports (GE MUSE, CLIF) are
converted by import formats (\file{medorama/src/formats/}). Removing a database takes out
everything it added. The requirements for tables are in \cref{chap:importing}.

\paragraph{Connected databases.} A \emph{connected database} (an OMOP CDM instance, or any SQL
database described by views) is read in place at each request, read-only, through three views of
canonical columns: persons, visits and entries (\file{medorama/src/sources/};
\cref{fig:connecting-sources}). Each person is linked to a catalog patient by the import rules
(patient ID, then MRN of the same authority), so a person in several sources is one patient with
one pseudonym. Cohort criteria are answered by every source that holds such data, and the
patient page shows every source's entries in the same columns. Everything leaves through the
same gate. How to connect one is described in \cref{chap:connecting}.

\paragraph{Recordings and indexing.} A recording (a WFDB record, an EDF file, a folder of
records, or a licensed format read by an add-on) is imported by path and never copied.
Background workers index it: the engine, or the add-on, reads the recording once and writes an
\emph{index} (channels, pieces, envelopes, vital signs, events) in the shared format of
\cref{chap:formats-reference}, into \code{index\_dir}. The viewers then open the recording
quickly from its index. Workers also derive series (for example a Holter's heart rate) and run
plug-in actions from the job queue (\file{jobs.rs}). The built-in \code{identify} action reads
the identifiers that a recording's own file states (for example an EDF header's patient field)
into the identifier tables, so that an administrator can check that the file belongs to its
patient.

\paragraph{Files read in place.} ECGs, DICOM studies, pictures and documents are served from
where they lie, one entry at a time, through \code{patients/\{id\}/file/\{item\}}. Medorama never
writes to them. In Privacy Mode, it de-identifies them as they are sent.

\section{Readers and data models}
\label{sec:architecture-readers}

Each file format has its own \emph{reader}, and every reader produces one of four \emph{data
models} (\cref{tab:architecture-models}).

\begin{xltabular}{\linewidth}{@{}P{3.2cm}P{6.0cm}L@{}}
\caption{The four data models.}
\label{tab:architecture-models}\\
\toprule
\tabhead Data model & \tabhead Produced by & \tabhead Used by \\
\midrule
\endfirsthead
\toprule
\tabhead Data model & \tabhead Produced by & \tabhead Used by \\
\midrule
\endhead
\bottomrule
\endfoot
Recording index & The engines' readers (WFDB, EDF and BDF) in Rust, C++ and Python; adapters
  for other formats & VitalStream; Medorama's indexing and cohorts \\
ECG record & The ECG viewer's readers (WFDB, GE MUSE XML, DICOM ECG, EDF, its JSON form) & The
  ECG viewer \\
DICOM objects & The PACS viewer's DICOM reader; \cmd{dicom-py}; the privacy gate's
  \file{dicom.rs} & The PACS viewer \\
Catalog tables (canonical columns) & \file{tables.rs}, import formats, connected databases'
  views & Patient pages, cohorts, the \ui{Databases} tab \\
\end{xltabular}

Past its reader, a format is like any other: the viewers, cohorts and the privacy gate know only
the data models. A new format is therefore one new reader. Where readers live and how to add one
is described in \cref{chap:readers}, and EDF's reading rules in \cref{chap:formats-reference}.

\paragraph{The engine, adapters and the toolbox.} \textbf{The engine} (\file{libs/engine}) holds
what is needed to read a recording quickly: the readers (WFDB records, multi-segment records,
folders of records read as one session, WFDB collections, EDF, EDF+, BDF and BDF+), the index,
the dataset queries that the viewer's API answers (windows, envelopes, vital signs, markers), and
a sample reader with which Medorama derives a Holter's heart rate. The Rust viewer is a thin HTTP
layer over it, and the C++ and Python viewers have their own engines that write the same index.
\textbf{Adapters} read formats that the engines do not read. An adapter is a program named
\code{<name>-adapter} in the adapter folder. The engine asks it to \code{describe} itself,
\code{detect} the files it reads, \code{index} a recording (writing the shared index format) and
\code{serve} samples over a pipe, and Medorama also asks it to \code{identify} a file's
identifiers. It may be written in any language (\cref{chap:addons}). \textbf{The toolbox}, which
is planned, will hold processing on waveforms (filters, beat detection, signal quality,
machine-learning inference), each tool available as a Medorama action, on the command line, from
Python and as tracks in the viewer.

\section{Extension points}
\label{sec:architecture-extension}

An institution adapts Medorama at a few documented places, without changing the rest
(\cref{tab:architecture-extension}).

\begin{xltabular}{\linewidth}{@{}P{3.4cm}P{7.2cm}L@{}}
\caption{Where Medorama is adapted, and the contract at each place.}
\label{tab:architecture-extension}\\
\toprule
\tabhead To adapt & \tabhead Where & \tabhead Described in \\
\midrule
\endfirsthead
\toprule
\tabhead To adapt & \tabhead Where & \tabhead Described in \\
\midrule
\endhead
\bottomrule
\endfoot
How users sign in & Configuration only, for an authenticating proxy (\code{sign\_in = "proxy"});
  or, in code, a new variant of \code{signin::Method} and its arm in \code{signin::identify}
  (\file{signin.rs}) & \cref{chap:sign-in} \\
A new recording format & A reader in the engine (\file{libs/engine/src/sources/mod.rs}), or an
  add-on in any language & \cref{chap:readers}, \cref{chap:addons} \\
A new ECG format in the ECG viewer & A reader registered in \file{viewers/ecg/readers.js} &
  \cref{chap:readers} \\
A new export to import & A TOML description in \code{formats\_dirs} (no code), or a converter in
  \file{medorama/src/formats/} & \cref{chap:importing} \\
An existing database, read in place & A connection of kind \code{omop}, or \code{sql} with views
  that give the canonical columns (\file{medorama/src/sources/}) & \cref{chap:connecting} \\
An analysis on recordings & A plug-in action: a folder with \file{action.toml} in
  \code{actions\_dirs} & \cref{chap:actions} \\
Scripts and other clients & The HTTP API (\file{medorama/openapi.yaml}) & \cref{chap:http-api} \\
\end{xltabular}

\section{Contracts between the parts}
\label{sec:architecture-contracts}

\begin{xltabular}{\linewidth}{@{}P{3.6cm}P{4.6cm}L@{}}
\caption{The contracts between the parts, and what checks them.}
\label{tab:architecture-contracts}\\
\toprule
\tabhead Contract & \tabhead Between & \tabhead Checked by \\
\midrule
\endfirsthead
\toprule
\tabhead Contract & \tabhead Between & \tabhead Checked by \\
\midrule
\endhead
\bottomrule
\endfoot
The index (folder layout, key, arrays; \cref{chap:formats-reference}) & The engines and add-ons,
  Medorama, the viewers sharing an index folder & \file{tests/parity/index\_parity.py} (byte for
  byte) \\
VitalStream's API, version 1 (\cref{chap:viewer-api}) & VitalStream's page and Medorama's
  gateway, and the three viewers & \file{tests/parity/viewer\_parity.py};
  \file{medorama/tests/gateway.rs} \\
The adapter protocol, version 1 (\cref{chap:addons}) & The engines and Medorama, and add-ons &
  The engines' adapter tests; \file{medorama/tests/sources.rs}; each add-on's conformance
  tests \\
EDF reading rules (\cref{chap:formats-reference}) & The three engines and the ECG viewer & The
  engines' EDF tests, the parity tests \\
WFDB collection (\cref{chap:formats-reference}) & The engines, and the tools that write long
  recordings as WFDB records & The viewers' tests, the parity tests \\
Viewer addresses (\code{\#src}, \code{\#list}, \code{\#path}) & Medorama and every viewer & The
  viewers' tests, Medorama's tests \\
Tables to import, connected databases' views (\cref{chap:importing}, \cref{chap:connecting}) &
  Source systems and Medorama & \file{medorama/tests/catalog.rs},
  \file{medorama/tests/connections.rs} \\
Medorama's HTTP API, version 1 (\cref{chap:http-api}) & Medorama and its page, scripts &
  \file{medorama/tests/openapi.rs} \\
Plug-in actions (\cref{chap:actions}) & Medorama and plug-ins & Medorama's tests \\
The catalog on SQLite and PostgreSQL & Medorama and its database & Every test, on both
  (\code{MEDORAMA\_TEST\_POSTGRES}) \\
\end{xltabular}

\section{Technology and dependencies}
\label{sec:architecture-technology}

\begin{itemize}
\item \textbf{Medorama} is written in Rust: axum and tokio for HTTP, rusqlite for SQLite
  (SQLCipher in encrypted builds), tokio-postgres with rustls for PostgreSQL, argon2 for
  passwords, and calamine and parquet for spreadsheets and Parquet files. Versions are pinned by
  \file{Cargo.lock} and checked by \cmd{cargo deny}.
\item \textbf{The web page and the ECG and PACS viewers} are plain HTML, CSS and JavaScript
  modules, with no build step and no third-party code, embedded in the server's binary.
\item \textbf{VitalStream} has three implementations: Rust (over the engine), C++20 (CMake) and
  Python. All three read WFDB through wfdb-x, a public WFDB library.
\item \textbf{The Python tools} \cmd{muse-py} and \cmd{dicom-py} build databases of GE MUSE ECGs
  and DICOM studies to import, and serve the ECG and PACS viewers alone.
\end{itemize}
Every third-party component and its license is listed in \cref{chap:licences}, and building is
described in \cref{chap:installation}.

\section{Code map}
\label{sec:architecture-code}

\Cref{tab:architecture-code} lists the main modules of the Medorama server
(\file{medorama/src/}). A guide for contributors is in the repository's
\file{docs/developing.md}.

\begin{xltabular}{\linewidth}{@{}P{5.6cm}L@{}}
\caption{The main modules of the Medorama server.}
\label{tab:architecture-code}\\
\toprule
\tabhead Module & \tabhead Responsibility \\
\midrule
\endfirsthead
\toprule
\tabhead Module & \tabhead Responsibility \\
\midrule
\endhead
\bottomrule
\endfoot
\file{server.rs} & Routes, the guard and privacy layers, files in Privacy Mode, background
  workers \\
\file{signin.rs} & Sign-in methods: password and proxy \\
\file{users.rs} & Accounts, roles, the PHI permission, sessions \\
\file{gate.rs}, \file{deid.rs}, \file{dicom.rs}, \file{pseudonyms.rs} & The privacy gate: the
  identifier dictionary and placeholders; dates, ages, WFDB, MUSE and EDF headers; DICOM;
  pseudonyms \\
\file{gateway.rs}, \file{viewer.rs} & The VitalStream gateway; the viewer's supervisor \\
\file{patient/}, \file{cohorts/}, \file{expr.rs}, \file{search.rs}, \file{query.rs} & The
  patient pages; cohorts, their criteria and exports; criteria written as text; search \\
\file{tables.rs}, \file{imports.rs}, \file{bulk.rs}, \file{typed.rs}, \file{formats/},
  \file{parquet\_rows.rs} & Importing databases; import formats \\
\file{linkage.rs}, \file{identity.rs}, \file{review.rs}, \file{ehr.rs} & Matching rows to
  patients and encounters; identity; the \ui{Review} tab; clinical entries \\
\file{databases.rs} & The \ui{Databases} tab and database exports \\
\file{sources/} & Connected databases: views, read-only pools, the person registry \\
\file{recordings.rs}, \file{identify.rs}, \file{derived.rs}, \file{categories.rs} & Recordings;
  identifiers stated in files; derived series; recording categories \\
\file{jobs.rs}, \file{actions.rs}, \file{addons.rs} & The job queue, plug-in actions, installed
  adapters \\
\file{audit.rs}, \file{db.rs}, \file{sql/}, \file{schema/}, \file{config.rs} & The sealed audit
  log; the database and its migrations; SQLite and PostgreSQL behind one interface; settings \\
\file{demo.rs} & The demo's invented patients \\
\end{xltabular}

\section{Status}
\label{sec:architecture-status}

Everything described above is built, except the toolbox and its results format. The remaining
work is listed item by item, with designs and tests, in the repository's \file{docs/todo.md}.
Next come per-study access lists, more of connected databases on the page, the toolbox and the
engine's streaming reader, and viewing recordings while they are being recorded. Later items are
server-side decoding of codecs that browsers do not decode (JPEG-LS, JPEG 2000, MPEG), DICOM
Holter recordings in VitalStream, free-text de-identification, FHIR, HL7 v2 and DICOMweb sources,
and ONNX inference.
